% Verbatim copy of /Users/alingalatan/workspaces/grothe/handoff/manuscript/legacy/ch56-coupled-tensor.tex (file dated 2026-10-25). Only this line was added.
% Consolidated source module (formerly Chapter 56). Owner: lower-bound/10-coupled-complements.tex
\section{Arbitrary coupled cubic tensors with unrestricted shared $A/B$}\label{ch:coupled-tensor}
\chapsummary{Arbitrary coupled remote tensors with unrestricted A/B under buffered variance/influence conditions. Actual nonlinear content is retained; small influence and fresh buffer are part of the theorem.}

\paragraph{Provenance and status (IDE note).} Full-context GPT, 20 October (\code{conjK-coupled-tensor-milestone-20261020.zip}, archived in
\pathref{gpt-reviews/coupled-tensor-20261020/}). IDE replay \code{python3 -I verify\_all.py} from an unrelated directory, regenerating the Gauss rule and the visible
Hankel table: all 12 continuum claims and both block corollaries PASS (150\,s; exact-branch margin $2.92\cdot10^{-5}$, weak-branch excess $8.33\cdot10^{-5}$).
IDE checks \st{P}: both branches of the content lemma (Cauchy--Schwarz with $\E|\ell|^2|q|^2\le e+h$; the exact integration by parts, $F_1'=t(tu-2)e^{-tu}$,
$\operatorname{Var}V=e^2+\kappa$, $\E V^2=2e^2+\kappa$), and the scalar step $\Lambda(q+\sqrt vU)\le F(v)+F'(v)e$ by concavity of $u\mapsto2\E|L+\sqrt vU+\sqrt u|$.
A byproduct worth recording: for $g\ge\sqrt{v_*}$ every remote cubic satisfies exact K, since $F'(g^2)\le0$.

\subsection{From independent drivers to arbitrary coupled cubic tensors}
\label{sec:coupled-tensor-milestone}

\paragraph{Status and scope.}
The analytic statements are \textbf{[P]}; the displayed numerical regions
are \textbf{[C]}, with a directed visible-table replay followed by exact
rational comparisons. Throughout,
\[
 a=699/1250,\quad b=179/1250,\quad p=2500/1891,\quad
 L=Z(a+b|Z|^2),
\]
\[
 \Lambda(G)=\mathbb E|L+G|+\mathbb E|L-G|-p\mathbb E|G|^2,\qquad
 S=\sup_{v\ge0}\{2\mathbb E|L+\sqrt v\,U|-pv\}.
\]
Let \(Z,U,X_1,\ldots,X_n\) be independent standard proper complex
Gaussians. The remote cubic \(q\in H_{2,1}(X_1,\ldots,X_n)\) is
arbitrary: it need not be a sum of independent Wick monomials, a
Gaussian factor, or independent of its shared linear drivers. Put
\[
 G=q+gU+(|Z|^2-1)\sum_j\alpha_jX_j+
                 \frac{Z^2}{\sqrt2}\sum_j\beta_j\bar X_j.
 \tag{CT.1}
\]
All \(\alpha_j,\beta_j\) are unrestricted complex numbers. The fresh
buffer \(U\) is independent of every original coordinate. Define
\[
 e=\mathbb E|q|^2,\quad
 \kappa=\mathbb E|q|^4-2e^2,\quad E_{AB}=\|\alpha\|^2+\|\beta\|^2,
\]
\[
 h=\sup_{\|\ell\|_2=1}
 \bigl(\|D_\ell q\|_2^2+\|\bar D_\ell q\|_2^2\bigr).
 \tag{CT.2}
\]
The supremum is over proper Gaussian holomorphic linear forms on the
remote space. Wirtinger derivatives are normalized by
\(\partial_{X_j}X_j=1\). Thus \(h\) is the largest eigenvalue of a
finite derivative Gram matrix, with the conjugations dictated by
\(D_\ell=\sum\bar u_j\partial_{X_j}\) and
\(\bar D_\ell=\sum u_j\partial_{\bar X_j}\) when
\(\ell=\sum u_jX_j\). In particular
\(D_\ell\bar q=\overline{\bar D_\ell q}\).
Always \(0\le h\le3e\) and \(0\le\kappa\le42e^2\).

\begin{theorem}[A finite-amplitude coupled-tensor region, \textbf{[P,C]}]
\label{thm:ct-influence}
For the full family (CT.1), in every finite dimension,
\[
 \boxed{
 g\ge\frac9{40},\quad e\le\frac1{25},\quad
 h\le\frac e{10}
 \quad\Longrightarrow\quad
 \Lambda(G)<S-\frac1{40000}-\frac{E_{AB}}{2000}.}
 \tag{CT.3}
\]
There is also a weak region nearer the optimal Gaussian buffer:
\[
 \boxed{
 g\ge\frac15,\quad e\le\frac1{25},\quad
 h\le\frac e{150}
 \quad\Longrightarrow\quad
 \Lambda(G)<S+\frac{89}{10^6}-\frac{E_{AB}}{2000}.}
 \tag{CT.4}
\]
The kurtosis restrictions in the original certificate are redundant:
Theorem~\ref{thm:sharp-influence-audit} gives $\kappa\le14eh$,
hence $\kappa\le7e^2/5$ in (CT.3) and $\kappa\le7e^2/75<e^2/10$
in (CT.4). The numerical proof below still checks the original
larger domain with $h$ and $\kappa$ separate; no recomputation is needed.

No amplitude or phase restriction on either shared sector is imposed.
The second statement uses the accepted remote \(10^{-5}\) theorem.
The first does not use that theorem, a sixth-moment bound, or any
unproved Gaussian-domination statement.
\end{theorem}

\begin{corollary}[Unequal arbitrary cubic blocks]
\label{cor:ct-blocks}
Suppose \(q=\sum_jq_j\), where the \(q_j\) are pure \(H_{2,1}\)
cubics on mutually independent remote Gaussian blocks. Each block
may contain an arbitrary internally coupled tensor. Write
\(e_j=\mathbb E|q_j|^2\), \(e=\sum e_j\le1/25\).
If \(\max e_j\le e/30\) and \(g\ge9/40\), then (CT.3) holds.
If \(\max e_j\le e/450\) and \(g\ge1/5\), then (CT.4) holds.
The A/B forms may use every coordinate of every block.
\end{corollary}
\begin{proof}
Circularity and independence give
\(\kappa=\sum_j\kappa_j\le42\sum_j e_j^2\le42e\max e_j\).
The derivative Gram matrix is block diagonal: derivatives from
distinct blocks are orthogonal, since each is a centered second-chaos
polynomial. Consequently \(h\le3\max e_j\).
For a block share \(1/30\), these are \(h\le e/10\) and
\(\kappa\le7e^2/5\). For \(1/450\), they imply
\(h\le e/150\), \(\kappa\le(7/75)e^2<e^2/10\).
\end{proof}

\paragraph{Influence-only form (reconciled).} The sharp bound
$\kappa\le14eh$ (Theorem~\ref{thm:sharp-influence-audit}) makes the
kurtosis condition in (CT.3) automatic under its original $h\le e/10$,
and the condition in (CT.4) automatic under $h\le e/150$.
Thus the exact branch allows $h\le e/10$, with its deficit $1/40000$.
The earlier $h\le9e/100$ certificate still retains the stronger deficit
$1/12500$. The displayed theorem now omits those redundant kurtosis assumptions; its underlying numerical certificate is unchanged. See
Corollary~\ref{cor:sharp-influence-region}.

\paragraph{An explicit mixed example with nonzero projected content.}
For sixty independent remote coordinates, put
\[
 B_j=\frac{H(X_j)+X_j^2\bar Y_j/\sqrt2}{\sqrt2},\qquad
 H(X)=X(|X|^2-2)/\sqrt2,\qquad
 q=\frac1{\sqrt{750}}\sum_{j=1}^{30}B_j .
\]
Each block has variance one and fourth moment \(26\).
Consequently \(e=1/25\), \(h\le e/10\), and \(\kappa=4e^2/5\).
Thus (CT.3) holds with \(g=9/40\) and arbitrary A/B forms on
\emph{all sixty coordinates}, including both \(X_j\) and \(Y_j\).
This example does not remove the nonlinear content by symmetry:
\[
 P_{1,0}(B_j|B_j|^2)=7X_j+Y_j .
\]
The moment and projection coefficients follow by conditioning on \(X_j\)
and using \(\mathbb E|X|^{2m}=m!\); the exact checker verifies the
finite arithmetic. Independence between the nonlinear coefficients
and all shared drivers is not available or assumed.

\subsubsection{Unconditional small-variance regions}
Without imposing any additional influence or kurtosis restriction,
the following peak-profile regions hold for every remote cubic \(q\)
and all complex A/B coefficients:
\[
\begin{array}{c|c|c}
 g\text{ at least}&e\text{ at most}&\text{conclusion}\\ \hline
 9/50&1/500&\Lambda<S+89/10^6-E_{AB}/2000\\
 37/200&9/4000&\Lambda<S+89/10^6-E_{AB}/2000\\
 1/5&7/2500&\Lambda<S+89/10^6-E_{AB}/2000\\ \hline
 9/40&49/10000&\Lambda<S-1/200000-E_{AB}/2000\\
 1/4&9/1000&\Lambda<S-1/200000-E_{AB}/2000\\
 11/40&7/500&\Lambda<S-1/200000-E_{AB}/2000\\
 3/10&19/1000&\Lambda<S-1/200000-E_{AB}/2000\\
 7/20&33/1000&\Lambda<S-1/200000-E_{AB}/2000\\
 2/5&49/1000&\Lambda<S-1/200000-E_{AB}/2000\\
 9/20&7/100&\Lambda<S-1/200000-E_{AB}/2000
\end{array}
\tag{CT.5}
\]
These are variance conditions on the actual remote cubic, not on its
Gaussian buffer or on the total A/B energy. All buffer intervals are
unbounded half-lines.

\subsubsection{A joint influence estimate, rather than a coordinate cover}
For \(k\ge0\), set
\[
 D(k)=\mathbb E[XX^*\cos(k\Re q)],\qquad
 C(k)=\mathbb E[XX^\top\cos(k\Re q)].
\]
For any unit Gaussian linear form \(\ell\), the finite Hermite
creation/annihilation identity is
\[
 \mathbb E|\ell|^2|q|^2
 =e+\|D_\ell q\|_2^2+\|\bar D_\ell q\|_2^2\le e+h.
 \tag{CT.6}
\]
For example, rotate \(\ell\) to \(X_1\), expand \(q\) in orthogonal
\(H_{\mu,\nu}\), and multiply by \(X_1\). The two orthogonal output
grades have squared norms with factors \(\mu_1+1\) and \(\nu_1\).
Their sum gives (CT.6), and \(\mu_1+\nu_1\le3\) gives \(h\le3e\).

The common phase rotation of \(X\) makes the \(q^2,\bar q^2\) terms
vanish in \(\mathbb E|\ell|^2(\Re q)^2\). Thus
\[
 \boxed{0\preceq I-D(k)\preceq k^2sI,\qquad
        \|C(k)\|_{\rm op}\le k^2s,\quad s=(e+h)/4.}
 \tag{CT.7}
\]
The first assertion uses \(1-\cos t\le t^2/2\).
For the second, write \(C=-\mathbb E[XX^\top w]\) with
\(w=1-\cos(k\Re q)\ge0\), and apply weighted Cauchy--Schwarz.
There is no entrywise summation and no dimension factor.

\subsubsection{Paying the actual nonlinear contents}
\label{sec:ct-content}
The accepted positive-kernel identities after visible averaging are
\[
 \mathbb E_Z[(|Z|^2-1)\overline{L+w}/|L+w|]
       =-\bar w k_A(|w|^2),\qquad
 \mathbb E_Z[Z^2\overline{L+w}/(\sqrt2|L+w|)]
       =-w k_B(|w|^2).
\]
Here \(k_A,k_B\) are positive completely monotone functions.
Gaussian buffering replaces each spectral atom by
\[
 \mathbb E[(w+gU)e^{-t|w+gU|^2}]
   =\frac{w}{(1+g^2t)^2}
           e^{-t|w|^2/(1+g^2t)}.
\]
The buffered kernels \(k_{j,g}\) are therefore completely monotone.
Write \(m_{j,1}=-k'_{j,g}(0)\).
Up to conjugation and sign, the content vectors are
\(P_{1,0}[qk_{j,g}(|q|^2)]\).

\begin{lemma}[Two bounds on the same projected content, \textbf{[P]}]
For a positive completely monotone \(k\), with \(m_1=-k'(0)<\infty\),
\[
 \boxed{\|P_{1,0}[qk(|q|^2)]\|_2^2\le m_1^2\,\Xi,\qquad
 \Xi=\min\{(e+h)(e^2+\kappa),\ h(6e^2+5\kappa)\}.}
 \tag{CT.8}
\]
\end{lemma}
\begin{proof}
Put \(V=|q|^2\), so \(\operatorname{Var}V=e^2+\kappa\).
First-chaos orthogonality permits subtracting \(k(e)q\).
For a unit \(\ell\), Cauchy--Schwarz and (CT.6) give
\[
 |\mathbb E[\bar\ell q(k(V)-k(e))]|^2
 \le (e+h)m_1^2\mathbb E(V-e)^2.
\]
This proves the first bound.

For the second use Gaussian integration by parts exactly:
\[
 \mathbb E[\bar\ell qk(V)]
 =\mathbb E[D_\ell q\{k(V)+Vk'(V)\}
              +q^2k'(V)D_\ell\bar q].
 \tag{CT.9}
\]
Both derivatives of \(q\) are centered second-chaos polynomials.
Set \(F_1(u)=k(u)+uk'(u)\).
For a spectral atom \(e^{-tu}\),
\[
 F_1'(u)=t(tu-2)e^{-tu},\qquad
 |(x-2)e^{-x}|\le2\quad(x\ge0).
\]
Consequently \(F_1\) is \(2m_1\)-Lipschitz and
\(\operatorname{Var}F_1(V)\le4m_1^2\operatorname{Var}V\);
the latter follows by taking an independent copy of \(V\).
Also \(\mathbb E|q^2k'(V)|^2\le m_1^2\mathbb E V^2\).
Subtract \(\mathbb EF_1(V)\) in (CT.9) and apply Cauchy--Schwarz
to the two-component pair. The result is
\[
 |\mathbb E[\bar\ell qk(V)]|^2
 \le h m_1^2\{4\operatorname{Var}V+\mathbb EV^2\}
 =h m_1^2(6e^2+5\kappa).
\]
Taking the supremum over unit \(\ell\) proves (CT.8).
All integrations by parts follow first for bounded spectral
truncations, then by dominated convergence. The required functions
and derivatives have polynomial growth and \(q\) is a Gaussian
polynomial.
\end{proof}

This estimate uses the actual shared drivers. In particular, its
second branch tends to zero with their largest directional influence.
It is not a bound obtained by substituting a cubic moment constraint
into each visible conditional slice.

\subsubsection{The global tangent and a two-by-two certificate}
Put \(T=|Z|^2\), \(\ell(T)=\sqrt T(a+bT)\), and define
\[
 R_{AA}=\mathbb E[(T-1)^2J_0(k\ell(T))],\quad
 R_{BB}=\mathbb E[T^2J_0(k\ell(T))/2],\quad
 R_{AB}=\mathbb E[T(T-1)J_2(k\ell(T))/\sqrt2].
\]
Let
\[
 W_g(k)=e^{-g^2k^2/4}-\sqrt\pi\,\frac{gk}{2}
                          \operatorname{erfc}(gk/2).
\]
At a fixed lower buffer \(g_0\), define the positive costs
\[
 A_0=p-\int_0^\infty W_{g_0}R_{AA}\,dk,\quad
 B_0=p-\int_0^\infty W_{g_0}R_{BB}\,dk,
\]
\[
 N_A=\int k^2W_{g_0}(R_{AA})_-\,dk,\quad
 N_B=\int k^2W_{g_0}(R_{BB})_-\,dk,\quad
 N_X=\int k^2W_{g_0}|R_{AB}|\,dk.
\]
The two kernel moments have the exact representations
\[
 m_{A,1}=\frac1{16}\int k^2e^{-g_0^2k^2/4}(-R_A(k))\,dk,\quad
 m_{B,1}=\frac1{16}\int k^2e^{-g_0^2k^2/4}R_B(k)\,dk,
 \tag{CT.10}
\]
where \(R_A=\mathbb E[(T-1)J_0(k\ell(T))]\le0\) and
\(R_B=\mathbb E[TJ_2(k\ell(T))/\sqrt2]\ge0\).
These formulas also fix the normalization of the contents:
the expansion \(J_1(x)/x=1/2-x^2/16+O(x^4)\) supplies the factor
\(1/16\).

The concavity of \(u\mapsto\mathbb E|\sqrt u+gU|\) gives a
global tangent at \(L+q\), with no size restriction on the A/B
perturbation. Remote parity \(X\mapsto-X\), with \(Z\) fixed,
identifies the two signs in the payoff. Combining its Fourier
representation with (CT.7) and (CT.8) gives, for
\(A=\|\alpha\|\), \(B=\|\beta\|\),
\[
 \Lambda(G)\le\Lambda(q+gU)
 +2\sqrt\Xi(m_{A,1}A+m_{B,1}B)
 -A_0A^2-B_0B^2+s(N_AA^2+2N_XAB+N_BB^2).
 \tag{CT.11}
\]
All bounds computed at \(g_0\) remain valid for \(g\ge g_0\).
Indeed \(W_g\) and the positive kernel moments decrease with \(g\).
The diagonal base integrals decrease by the physical-space identity
\[
 K_g(u)=\frac1{\sqrt\pi}\int_0^{1/g^2}
           t^{-1/2}\sqrt{1-g^2t}\,e^{-tu}\,dt,
\]
not by an unjustified sign assumption on their Hankel integrands.

Suppose now \(e\le E\), \(h\le re\), \(\kappa\le de^2\).
Set \(s_0=(1+r)E/4\), \(\tau=1/2000\), and use directed upper
or lower constants as appropriate in
\[
 P=\begin{pmatrix}A_0&0\\0&B_0\end{pmatrix}
 -s_0\begin{pmatrix}N_A&N_X\\N_X&N_B\end{pmatrix}-\tau I.
 \tag{CT.12}
\]
The replacement \(s\le s_0\) is made on the nonnegative norm pair
\((A,B)\); no Loewner monotonicity of the matrix of \(N\)'s is
asserted. If \(P\) is positive definite, completing the square yields
\[
 \Lambda(G)\le\Lambda(q+gU)+K_*e^3-\tau E_{AB},
\]
\[
 K_*=
 \frac{\min\{(1+r)(1+d),\,r(6+5d)\}}{176}
 u^\top P^{-1}u,\qquad
 u=\sqrt{176}(m_{A,1},m_{B,1})^\top.
 \tag{CT.13}
\]
The numerical output stores upper bounds for the two entries of \(u\).
The inverse of \(P\) has nonnegative entries, so these substitutions
are valid. Exact fractions check its principal minors.

\subsubsection{Closing the scalar baseline}
\paragraph{Exact branch: no remote comparison hypothesis.}
The visible Gaussian scale mixture \(L\overset d=\sqrt V\,W\)
implies that
\[
 f_v(u)=2\mathbb E|L+\sqrt v\,U+\sqrt u|
\]
is concave in \(u\). Write \(F(v)=f_v(0)-pv\).
Its derivative satisfies \(f_v'(0)=F'(v)+p\).
The tangent at zero therefore proves for every remote law of variance \(e\)
\[
 \Lambda(q+\sqrt v\,U)\le F(v)+F'(v)e.
 \tag{CT.14}
\]
Moreover \(F\) is concave in \(v\). If the directed computation gives
\(F'(g_0^2)<0\), then for every \(g\ge g_0\) the right side of
(CT.14) is at most \(F(g_0^2)+F'(g_0^2)e\).
Since \(e^3\le E^2e\), it is enough to check the two endpoints
\[
 \max\{F(g_0^2),\,
 F(g_0^2)+F'(g_0^2)E+K_*E^3\}<S_{\rm lo}-\delta.
 \tag{CT.15}
\]
This is an exact comparison of rational interval endpoints.
The benchmark lower bound is computed within this package:
\[
 S_{\rm lo}=F(1/25)_{\rm lower}>1.3739533908.
 \tag{CT.16}
\]

\paragraph{Weak branch.}
The accepted remote-cubic theorem gives
\(\Lambda(q+gU)<S+10^{-5}\), in every dimension and for every
fresh buffer. Therefore it suffices to check
\[
 10^{-5}+K_*E^3<89/10^6.
 \tag{CT.17}
\]
This use of the remote theorem is restricted to the genuinely
remote baseline, before either visible shared sector is inserted.

\paragraph{Numerical proof and continuum coverage.}
The checker verifies (CT.12), (CT.15), or (CT.17) for all twelve
displayed claims. For (CT.3) the exact interval comparison leaves
a margin greater than \(2.9237\cdot10^{-5}\), sufficient for
\(1/40000\). For (CT.4) the certified excess is less than
\(8.3302\cdot10^{-5}\), sufficient for \(89/10^6\).
Decimals here are reports, not proof inputs. The monotonicity and
affine endpoint arguments above cover every intermediate variance
and the entire buffer half-line.

The universal bounds used in (CT.5) are \(r=3,d=42\).
The fourth-moment estimate \(\mathbb E|q|^4\le44e^2\) follows,
for \(q=I_{2,1}(f)\), by expanding \(q^2\) into grades
\((4,2),(3,1),(2,0)\) with coefficient bounds \(1,4,4\).
Their squared isometry factors sum to \(48+96+32=176\), while
\(e=2\|f\|^2\). This is the established finite product calculation;
no dimension-dependent hypercontractivity constant is introduced.

\subsubsection{Numerical method, dependencies, and limitations}
The inherited audited generator supplies outward enclosures of the
nine visible Hankel functions at 8,400 Gauss nodes, with analytic
quadrature and visible tails. No remote-driver interpolation or
parameter-box cover is used in this package.
For the nonsmooth negative parts and absolute value defining the
three \(N\)'s, the calculation integrates the analytic functions
\(R\) and \(R^2\) and applies Cauchy--Schwarz with
\(d\mu=k^2W_g\,dk\), \(\mu(\mathbb R_+)=\sqrt\pi/(2g^3)\).
It does not apply an analytic Gauss remainder to \(R_-\).

The extra baseline integrals use
\[
 \frac{F(v)+pv}{2}
 =\int_0^\infty\frac{1-e^{-vk^2/4}R_0(k)}{k^2}\,dk,\qquad
 F'(v)+p=\frac12\int_0^\infty e^{-vk^2/4}R_0(k)\,dk.
\]
The integrand at zero is evaluated through its removable-singularity
formula. At \(k=70\), the first tail is \(1/70\) plus an enclosed
Gaussian remainder; the second has a Gaussian tail bound.
Every final comparison uses exact fractions of binary64 interval
endpoints. The code disables fast math and fused contraction.

Inputs are the visible Gaussian-mixture and positive-kernel
identities, elementary Gaussian integration by parts and Hermite
orthogonality, and the established fourth-moment inequality.
Only the weak claims additionally use the accepted remote
\(10^{-5}\) theorem and its moment inputs.

These theorems cover genuinely coupled tensors, but require a fresh
buffer and either small remote variance or explicit influence and
fourth-moment conditions. They do not cover arbitrary even or
diagonal additions after A/B. (\emph{Later:} \autoref{ch:all-sector-audit-extension} removes the even/diagonal gap in six buffered coupled-cubic regions, not at zero buffer.) The earlier full-sector local theorem
is therefore not superseded. Profile uniformity, a complete row
cover, and the remaining general-tensor regions are open. No new
Grothendieck endpoint is claimed.
